%%%%%%%%%%%%%%%%%%%%%%%%%%%%%%%%%%%%%%%%%%%%%%%%%%%%%%%%%%%%%
% ==================== 可选：模型的分析与检验 ====================
% 本文件默认不在 main.tex 中启用。
% 仅当全文采用同一模型或同类模型，且需要统一比较检验指标、关键参数或误差传播时启用。
% 启用后必须删除各问题中重复的“模型检验与灵敏度分析”小节，避免重复论证。
% ✓ 自查清单：
%   □ 是否说明被统一检验的模型、指标和判定标准？
%   □ 参数选择及变化范围是否有现实或统计依据？
%   □ 是否分析输入扰动、假设变化、误差来源与传播？
%   □ 所有指标和结论是否来自真实计算记录？

\section{模型的分析与检验}

本文对【需要统一分析的模型或同类模型】进行集中检验。
根据各模型承担的论文任务，分别选取【量化指标】评价【拟合、预测、排序、约束可行性或求解稳定性】，并以
【理论条件、统计阈值、基准方法或题目要求】作为判定标准。真实检验结果见【图表或结果位置】，由此可得【总体检验结论】。

\subsection{灵敏度分析}

选取具有现实不确定性且可能改变主要结论的参数【参数名称】进行灵敏度分析，其基准值和变化范围分别为
【基准值】和【有依据的变化范围】，选择依据是【数据误差、现实波动或参数来源】。
在保持【控制变量】不变的条件下重新求解模型，记录【目标值、关键状态或决策方案】的变化。
结果表明【真实变化规律】，模型在【稳定区间】内保持【结论或方案】不变；当参数达到【临界值】时，
【数值或决策结构】发生【真实变化】，说明【敏感性结论及现实含义】。

\subsection{误差与稳健性分析}

模型误差主要来自【测量、抽样、参数估计、模型简化或数值计算等真实来源】，其可能通过
【变量、方程或求解步骤】传递至【核心输出】。采用【误差指标、输入扰动、假设替换、重复实验或压力测试】
进行分析，得到【真实误差或稳定性结果】。与【允许误差、置信区间或基准方案】比较后可知，
模型在【适用范围】内具有【稳健性结论】；在【边界条件】下仍存在【需要说明的风险或局限】。
